\section{Why Single-Fault Damage Is Not Protection Value}\label{sec:allocation}
This section states what can and cannot be expected of a static ranking. Its content is standard reliability theory applied to bit protection.

The isolated damage of bit $b$, $I_b=\risk(\{b\})-\risk(\varnothing)$, is the risk after one flip minus the clean risk. Protecting a bit, however, removes its effect in a network that may already be faulty. For an unprotected bit $b$, let $G$ be the random set of faults on the other unprotected bits, and define the conditional vulnerability $V_b(p,\protset)=\E_G[\risk(G\cup\{b\})-\risk(G)]$. Because bit $b$ faults independently of $G$,
\begin{equation}\label{eq:protection-benefit}
J_p(\protset)-J_p(\protset\cup\{b\})=p\,V_b(p,\protset).
\end{equation}
The left side is the expected risk reduction from protecting $b$. This is the improvement potential of reliability theory, the product of the failure probability and Birnbaum's importance~\cite{birnbaum1969importance,rausand2004system}. $V_b$ equals $I_b$ when the risk is additive over faults or when no other fault is present, and in general depends on both $p$ and the current mask.

Two small examples show the consequences. With $z_i=1$ denoting a flip, let
\begin{equation}\label{eq:three-bit-risk}
\risk(z_1,z_2,z_3)=1.2z_1+z_2+0.1z_3+3z_2z_3.
\end{equation}
Isolated damage ranks bit one first. Protecting it leaves expected risk $1.1p+3p^2$, whereas protecting bit two leaves $1.3p$, so the better single bit to protect switches at $p=1/15$. No order of these bits that ignores $p$ is best at every $p$.

A ranking also forces masks to be nested across budgets, and optimal masks need not be. For $\risk=z_1+1.05z_2+1.05z_3+z_1z_2+z_1z_3$ at $p=0.1$, the best single protected bit is bit one (expected risk $0.210$ against $0.215$), but the best protected pair is bits two and three ($0.100$ against $0.105$).

These examples prove only that no guarantee exists. They do not say that a static ranking performs badly in a real network, nor how large the loss from using one is. A fixed mask can be optimal for a fixed condition, and a static score may itself be built from multi-fault measurements. How much is lost in practice is an empirical question, which the remaining sections address by measuring task quality on the full BER-by-budget grid.
